% s4-import.tex -- Section 4 of the reader's companion (drafter: cqft, R3; edited in R4).
% Retitled and rewritten for the revised edition of 7 November 2026: the uniform local
% moment is a theorem of the record (Appendix D); A8(vi) is retired; A8(iv)-T is disclosed.
% Status sentence: ugjs/PLAN.md Part 8, F-B with the A8(iv)-T variant, verbatim.
\section{The uniform local moment, from the printed clustering theorem}\label{sec:import}

\emph{The uniform local moment is a theorem of the record: restrict the cutoff to a box of
GJS squares; Theorem~1.1.7 bounds the difference, Appendix~C the box.}

Of the record's eight imports \rec{\S2.4}{9}, A1--A7 quote printed statements (table
below). A8 packages the input of Glimm--Jaffe--Spencer \cite{GJS} for the
interacting Euclidean measures $dq_h=Z(h)^{-1}e^{-\mathcal U_h}d\mu_0$, $h\in\CE$, of A3
($\mathcal U_h$ the Euclidean interaction, $d\mu_0$ the free Euclidean measure): the
cutoff classes, the observables, the printed norms $\|\cdot\|_{\rm GJS}$, the dilation, and
one theorem, CE2; its items (v)--(vii$'$) are provenance, used by no proof
\rec{A8, \emph{Used in}}{19}.

Theorem~12.2 is proved from printed statements only, with one GJS theorem imported:
Theorem~1.1.7, which is also used for sharp-time observables across strips transverse to
the time-zero line, a case its printed statement covers and its printed proof does not
reach (listed separately as A8(iv)-T). The uniform local moment is a theorem of the record
(Appendix~D). The human readings of C5 items 1--4 and 10--12, with item 9(c), are open
(\S\ref{sec:record}). The imported theorem, as the record transcribes it:

\begin{recordbox}{Source Fact 2.2 of the record: CE2 $=$ GJS Theorem 1.1.7 (p.\,15)}
For localized monomials $Q_1,Q_2$ with disjoint localization squares separated by a strip
of width $d$ bounded by two parallel lines (any orientation),
\[
  \Bigl|\int Q_1Q_2\,dq_h-\int Q_1dq_h\int Q_2dq_h\Bigr|
  \le C_7e^{-m_1d}\|Q_1\|_{\rm GJS}\|Q_2\|_{\rm GJS} ,
\]
with $m_1>0$ and $C_7$ \emph{independent of $h$, $Q_1$, $Q_2$} (so stated in the source).
\end{recordbox}

\emph{The comparison.} The record's \S\S5 and~7 need moments of sharp-time observables
under the slab measures $dq_{h_{T,g}}$, bounded uniformly in $T$. Appendix~D derives them
with the device of GJS's own proof of their (1.1.11) \rec{Appendix~D}{53}, in three steps
(Fig.~\ref{fig:quantifier}); $\langle F\rangle_h:=\int F\,dq_h$.
(1)~Restrict $h\in\CE$ to a box $B$ of GJS squares (unit squares centred at integer
points) containing the squares of the observable $F$, $h^B:=h\mathbf 1_B$. Along
$h_\alpha:=h^B+\alpha(h-h^B)\in\CE$, $\alpha\in[0,1]$, the difference
$\langle F\rangle_h-\langle F\rangle_{h^B}$ is minus the $\alpha$-integral of the
covariance, under $dq_{h_\alpha}$, of $F$ with the interaction $\mathcal U_{h-h^B}$ outside
the box \rec{Lemma~D.1}{79}.
(2)~That interaction is a sum over the squares $\Delta$ outside $B$ charged by $h$; a strip
of positive width $w_\Delta$ separates each from the squares of $F$, and Theorem~1.1.7 bounds its
contribution by $C_7c_U\|F\|_{\rm GJS}e^{-m_1w_\Delta}$, with $c_U$ a bound on the GJS norm
of the interaction in one square \rec{Lemma~D.2}{84}.
(3)~The box term needs no clustering: Lemma~C.2 bounds $Z(h^B)$ from below, and
$\|e^{-\mathcal U_{h^B}}\|_{L^2(d\mu_0)}$ is at most $S_2(B)$, the supremum of that norm over
the cutoffs supported in the closed box: finite by Theorem~C.5, not evaluated, and
dependent on the shape of $B$ only \rec{Proposition~N.15}{62}. Two choices of the box give
two theorems, the second \Ttext{resting on A8(iv)-T}:
\begin{checkpoint}[Checkpoint: Theorems D.3 and D.4]
Assume the standing assumption of \rec{\S D.c.1}{82}. For $g\in\CH$, a bounded interval $I$ of positive length, $s\in\R$ and a finite real
linear combination $F$ of cylinder monomials supported in $I$, for every $T>0$,
\begin{equation*}
 \Bigl|\int F(s)\,dq_{h_{T,g}}\Bigr|\le A\,\|F\|_{L^2(d\phi_0)}+B'\,\|F(s)\|_{\rm GJS} ,
 \tag*{\rec{Theorem~D.3}{90}}
\end{equation*}
and for every $h\in\CE$, $a\in\R$ and real measurable $h'$ with $\|h'\|_\infty\le1$ and
$\supp h'\subseteq[a,a+2]$,
\begin{equation*}
 \int V(h')(0)^2\,dq_h\le M_B^2 .
 \tag*{\rec{Theorem~D.4}{95}}
\end{equation*}
\end{checkpoint}
Here $F(s)$ is $F$ at time $s$, $V(h')(0)$ the local interaction at time $0$,
$A=e^{3\lambda|a_0|N_B}S_2(B)$, $B'=2C_7c_UN_ge^{-m_1}(1-e^{-m_1})^{-1}$, with $N_g\le N_B$
the column counts of $\supp g$ and of the box; $M_B$ is in \S\ref{sec:euclidean}.

For Theorem~D.3 the box is three rows of squares around the time $s$, across every column
from the leftmost to the rightmost one meeting $\supp g\cup I$; every charged square outside lies
above or below, so every strip is parallel to the time-zero line \rec{Remark~N.39}{92}.
The constants then depend on $g$, which is all the record's \S7 needs: its limits
$T\to\infty$ are taken at fixed $g$. For Theorem~D.4 the observable lies in three
adjacent squares of row $0$, and $B_{3\times5}$ adds a row above and below and a column on
each side. The squares outside lie in rings, the $j$-th of $8j+12$ squares: its top and
bottom rows, class (P), are separated from the observable by a parallel strip of width
$j$; its side columns, class (T), by a transverse strip of width $j$, and in rows
$-1,0,1$ by no parallel strip of positive width at all. With parallel strips alone the
sum would grow with the number of columns that $h$ charges, which is not bounded uniformly
in $h\in\CE$; the transverse strips make $M_B$ one number for every $h$
\recpp{\S D.e.4}{96--97}.

\emph{A8(iv)-T.} The record discloses: ``GJS Theorem 1.1.7 is printed for localized
monomials separated by a strip of any orientation (p.~594). Its printed proof rotates the
strip into the time direction (p.~597) and so reaches sharp-time monomials only across
strips parallel to the time-zero line. Lemma~7.4 uses only that case; Appendix~D also uses
the transverse case, listed separately as A8(iv)-T'' \rec{A8(iv-T)}{15}. In the record's
uses the second monomial is integrated over time and the strip is bounded by two lines of
constant $x^1$. The case enters Theorem~D.4, hence \Ttext{Theorem~5.1 and every statement
that consumes $M$}, from Lemma~10.2 on; it does not enter the record's \S7, whose
Theorem~D.3 uses parallel strips only. It is listed for transparency, not as another kind
of import: the printed statement is used at its printed scope, as every item of A8 is. It is not a narrower form of the retired A8(vi), a moment bound on the slab
class that stated no clustering; neither statement is known to imply the other
\rec{A8(iv-T)}{15}. The record does not prove it \rec{Appendix~D}{53}; closing it is
Phase~C of \texttt{ugjs/PLAN.md}, optional and not scheduled (\texttt{docs/imports.md}).
Nor does the record re-derive the cluster expansion behind Theorem~1.1.7, whose printed
proof takes its constant from Theorem~1.1.8 at the same $h$ (p.~597)
\rec{A8, Scope~(c)}{18}.

\emph{Appendix~C, and the history.} Theorem~12.2 now consumes Appendix~C: Lemmas~C.2 and
C.4 and Theorem~C.5 (with Lemma~C.3) enter Theorems~D.3--D.4 \rec{Appendix~C}{45}, with
inputs the record repaired on 7~November 2026 (\S\ref{sec:record}); its audit of the source
mechanism is kept, Proposition~C.1 being ``deliberately limited to the audited
mechanisms'' \rec{Proposition~C.1}{47}. From 16~August to 30~October 2026 the record
carried the cutoff-uniform moment estimate $\UGJS$, A8(vi), as a named hypothesis of
Theorem~12.2; on 31~October it read that item off one printed theorem, GJS
Theorem~1.1.11 (a re-sourcing, not a proof). Since 7~November 2026 A8(vi) is retired:
Theorems~D.3--D.4 supply what the record's \S\S5 and 7 need, and no proof uses it
\rec{A8, Scope~(a)}{18}.

\emph{The seven literal imports} \rec{\S2.4}{9}:

\begin{center}
\footnotesize
\setlength{\tabcolsep}{4pt}
\begin{tabular}{@{}>{\raggedright\arraybackslash}p{0.04\linewidth}
                >{\raggedright\arraybackslash}p{0.55\linewidth}
                >{\raggedright\arraybackslash}p{0.35\linewidth}@{}}
\toprule
& import: what it supplies & source pinpoint\\
\midrule
A1 & \emph{cutoff construction:} the ground state $\Om_g$ of $H(g)$; the closure of $K(g)$
``has $E(g)=\inf\spec K(g)$ as a \emph{simple} eigenvalue''; the Wick powers in $V(h)$ &
\cite[p.~588]{GJS}; \cite[Thms.~2.2.1, 2.3.1]{GJ2}; \cite{GJ71}; \cite[Thm.~5.3]{Rosen70};
\cite[pp.~10--12, 15--17, 134--138]{SimonBook}\\
A2 & \emph{Feynman--Kac--Nelson transfer:} the slab objects $h_{T,g}$, $\psi_s$, $c_g$
and the Euclidean $\leftrightarrow$ Hamiltonian passage &
\cite[Thm.~III.6, Cor.~V.14, Thm.~V.15]{SimonBook}; \cite[Thms.~II.14, II.16]{GRS75}\\
A3 & \emph{finite-volume stability:} $\mathcal U_h$, built from Simon's Wick powers, in
every $L^p$, and $e^{-\varrho\,\mathcal U_h}$ integrable ($\varrho>0$): the measures
$dq_h$ are well defined &
\cite[pp.~135--137; Thm.~V.2(a), p.~137; Thm.~V.7(a), p.~152]{SimonBook}\\
A4 & \emph{free hypercontractivity:} the $L^p$ smoothing of the pairing bound (the
record's \S8) and the moment comparison on the chaoses (Appendices~C, D) &
\cite[Thms.~I.17, I.22, III.17]{SimonBook}\\
A5 & \emph{domain smoothing; common core:} $\Om_g\in D(V(h))$, and a common core for all
$K(g_s)$ along an affine path & (i) \cite[Ch.~7, Thm.~7.4]{GJExp}; (ii)
\cite[Thm.~V.12]{SimonBook}, \cite[Thm.~5.3]{Rosen70}, \cite{GJ2}\\
A6 & \emph{free propagation; patching; the dynamics $\alpha$:} (i) speed one, the
Trotter formula; (ii) (P): $\alpha_t(\A(O))\subseteq\A(O_{|t|})$,
$\alpha_t(A)=\beta^g_t(A)$ if $g\equiv1$ on $O_{|t|}$ ($A\in\A(O)$); (iii) locality:
$\A(O)\subseteq\A(O')'$ for disjoint $O$, $O'$; the cutoff cone is Proposition~2.7 &
\cite[\S III]{GJ1}; \cite[\S\S3.5--3.6]{GJ2}; \cite[\S4.1 and Ch.~4]{GJExp};
(iii) \cite[p.~397]{GJ2}\\
A7 & \emph{covariance data:} translation covariance and parity & \cite[p.~398]{GJ2}\\
\bottomrule
\end{tabular}
\end{center}

\noindent\emph{In the record:} \S2.4 with A1--A8, Source Fact~2.2, item (iv-T) and the
Scope paragraph, pp.~9--19; Appendix~C (Proposition~C.1, Lemma~C.2, Theorem~C.5),
pp.~45--51; Appendix~D (Lemmas~D.1--D.2, Proposition~N.15, Remark~N.39, Theorems~D.3--D.4,
\S D.e.4), pp.~53--97.
